\documentclass[10pt,a4paper]{amsart}
\usepackage{amsmath,amssymb,amsthm,graphicx}
% Retained minimal preamble for the Stacks Project chapter sites-cohomology.tex (Cohomology on Sites).
% Verbatim excerpt of the chapter preamble: Xy-pic setup (native preamble lines 14-18) and the native
% theorem declarations with the mathematical macros (native preamble lines 154-217). Copied mechanically;
% nothing added, removed or redefined. The wrapper class and the amsmath/amssymb/amsthm/graphicx base
% packages are supplied by the assembler, so the class/externaldocument part of the native preamble is omitted.
% For commutative diagrams we use Xy-pic
\usepackage[all]{xy}

% We use 2cell for 2-commutative diagrams.
\xyoption{2cell}
% Theorem environments.
%
\theoremstyle{plain}
\newtheorem{theorem}[subsection]{Theorem}
\newtheorem{proposition}[subsection]{Proposition}
\newtheorem{lemma}[subsection]{Lemma}

\theoremstyle{definition}
\newtheorem{definition}[subsection]{Definition}
\newtheorem{example}[subsection]{Example}
\newtheorem{exercise}[subsection]{Exercise}
\newtheorem{situation}[subsection]{Situation}

\theoremstyle{remark}
\newtheorem{remark}[subsection]{Remark}
\newtheorem{remarks}[subsection]{Remarks}

\numberwithin{equation}{subsection}

% Macros
%
\def\lim{\mathop{\mathrm{lim}}\nolimits}
\def\colim{\mathop{\mathrm{colim}}\nolimits}
\def\Spec{\mathop{\mathrm{Spec}}}
\def\Hom{\mathop{\mathrm{Hom}}\nolimits}
\def\Ext{\mathop{\mathrm{Ext}}\nolimits}
\def\SheafHom{\mathop{\mathcal{H}\!\mathit{om}}\nolimits}
\def\SheafExt{\mathop{\mathcal{E}\!\mathit{xt}}\nolimits}
\def\Sch{\mathit{Sch}}
\def\Mor{\mathop{\mathrm{Mor}}\nolimits}
\def\Ob{\mathop{\mathrm{Ob}}\nolimits}
\def\Sh{\mathop{\mathit{Sh}}\nolimits}
\def\NL{\mathop{N\!L}\nolimits}
\def\CH{\mathop{\mathrm{CH}}\nolimits}
\def\proetale{{pro\text{-}\acute{e}tale}}
\def\etale{{\acute{e}tale}}
\def\QCoh{\mathit{QCoh}}
\def\Ker{\mathop{\mathrm{Ker}}}
\def\Im{\mathop{\mathrm{Im}}}
\def\Coker{\mathop{\mathrm{Coker}}}
\def\Coim{\mathop{\mathrm{Coim}}}

% Boxtimes
%
\DeclareMathSymbol{\boxtimes}{\mathbin}{AMSa}{"02}

%
% Macros for moduli stacks/spaces
%
\def\QCohstack{\mathcal{QC}\!\mathit{oh}}
\def\Cohstack{\mathcal{C}\!\mathit{oh}}
\def\Spacesstack{\mathcal{S}\!\mathit{paces}}
\def\Quotfunctor{\mathrm{Quot}}
\def\Hilbfunctor{\mathrm{Hilb}}
\def\Curvesstack{\mathcal{C}\!\mathit{urves}}
\def\Polarizedstack{\mathcal{P}\!\mathit{olarized}}
\def\Complexesstack{\mathcal{C}\!\mathit{omplexes}}
% \Pic is the operator that assigns to X its picard group, usage \Pic(X)
% \Picardstack_{X/B} denotes the Picard stack of X over B
% \Picardfunctor_{X/B} denotes the Picard functor of X over B
\def\Pic{\mathop{\mathrm{Pic}}\nolimits}
\def\Picardstack{\mathcal{P}\!\mathit{ic}}
\def\Picardfunctor{\mathrm{Pic}}
\def\Deformationcategory{\mathcal{D}\!\mathit{ef}}
\usepackage{cleveref}
\crefname{originallemma}{original lemma}{original lemmas}
\Crefname{originallemma}{Original lemma}{Original lemmas}
\makeatletter
\newlabel{derived-lemma-higher-derived-functors}{{015B}{source}{}{sourcepointer-1}{}}
\newlabel{derived-lemma-higher-derived-functors@cref}{{[originallemma][0][]015B}{[1][0][]source}{}{}{}}
\makeatother
\crefname{originallemma}{original lemma}{original lemmas}
\Crefname{originallemma}{Original lemma}{Original lemmas}
\makeatletter
\newlabel{homology-lemma-uniqueness-universal-delta-functor}{{010U}{source}{}{sourcepointer-2}{}}
\newlabel{homology-lemma-uniqueness-universal-delta-functor@cref}{{[originallemma][0][]010U}{[1][0][]source}{}{}{}}
\makeatother
\crefname{originallemma}{original lemma}{original lemmas}
\Crefname{originallemma}{Original lemma}{Original lemmas}
\makeatletter
\newlabel{homology-lemma-efface-implies-universal}{{010T}{source}{}{sourcepointer-3}{}}
\newlabel{homology-lemma-efface-implies-universal@cref}{{[originallemma][0][]010T}{[1][0][]source}{}{}{}}
\makeatother
\crefname{originallemma}{original lemma}{original lemmas}
\Crefname{originallemma}{Original lemma}{Original lemmas}
\makeatletter
\newlabel{lemma-cech-vanish-collection}{{03F9}{source}{}{sourcepointer-4}{}}
\newlabel{lemma-cech-vanish-collection@cref}{{[originallemma][0][]03F9}{[1][0][]source}{}{}{}}
\makeatother
\crefname{originallemma}{original lemma}{original lemmas}
\Crefname{originallemma}{Original lemma}{Original lemmas}
\makeatletter
\newlabel{lemma-injective-module-trivial-cech}{{03FC}{source}{}{sourcepointer-5}{}}
\newlabel{lemma-injective-module-trivial-cech@cref}{{[originallemma][0][]03FC}{[1][0][]source}{}{}{}}
\makeatother
\crefname{originallemma}{original lemma}{original lemmas}
\Crefname{originallemma}{Original lemma}{Original lemmas}
\makeatletter
\newlabel{sites-lemma-topos-good-site}{{03CI}{source}{}{sourcepointer-6}{}}
\newlabel{sites-lemma-topos-good-site@cref}{{[originallemma][0][]03CI}{[1][0][]source}{}{}{}}
\makeatother
\title{Cohomology of a module over a sheaf of rings agrees with cohomology of the underlying abelian sheaf}
\author{The Stacks Project Authors}
\date{stacks-project:sites-cohomology.tex:local-snapshot-20261011; GNU FDL 1.2 or later}
\begin{document}
\maketitle
\noindent\textit{Editorial note.} All delivered mathematical text is the author's own, copied from the pinned Stacks Project chapter file sites-cohomology.tex; no statement or proof step was rewritten, reordered or removed.\par\medskip
\noindent\textit{Editorial note.} The proof is the author's complete argument, including the discussion of the case where the site has no final object.\par\medskip
\noindent\textit{Editorial note.} Every labelled reference occurring inside the counted spans is carried in assembly.external\_references as a non-load-bearing navigation pointer (same-chapter ones at their definition site, cross-chapter ones at the invocation inside the counted span), so no load-bearing mathematics is replaced by a stub. No bibliography citation occurs inside the counted spans.\par\medskip
\noindent\textit{Editorial note.} The checked chapter list for same-name collisions is the set of stacks chapters staged in this repository: chow, cohomology, divisors, etale-cohomology, flat, more-algebra, more-morphisms, morphisms, spaces-more-morphisms, varieties together with this card's chapter sites-cohomology.\par\medskip
\noindent\textit{Editorial note.} The preamble is the chapter's own retained minimal preamble (sources/stacks-sites-cohomology/retained-preamble.tex), a verbatim excerpt of the native preamble: the Xy-pic setup and the native theorem declarations with the mathematical macros.\par\medskip
\noindent\textit{Allowed standard prerequisites (editorial).}\begin{itemize}
\item Cohomology in Ab(C) and in Mod(O) of a ringed site. Same chapter, the two cohomology functors on abelian sheaves and on modules over a sheaf of rings, computed from injective resolutions in each category. Supplies the two sides of the asserted equality and the resolutions with which they are computed.
\item Vanishing criterion for totally acyclic sheaves over a cofinal system. Same chapter, the companion lemma giving the vanishing of cohomology from the vanishing of Cech cohomology over a cofinal system of coverings. Supplies the criterion applied to an injective module in the proof.
\end{itemize}\medskip
\ifdefined\hypertarget\hypertarget{sourcepointer-1}{}\fi
\noindent\textit{Source reference.} Original lemma 015B: sources/stacks-sites-cohomology/pinned-original-source.tex, source lines 1886--1886. Navigation pointer invoked at proof line 1886; the counted span contains the pointer only and the cited result is not reproduced inside it, so the pointer is registered as a non-load-bearing navigation stub in the established form for Stacks material. No bibliography citation is involved.\par\medskip
\ifdefined\hypertarget\hypertarget{sourcepointer-2}{}\fi
\noindent\textit{Source reference.} Original lemma 010U: sources/stacks-sites-cohomology/pinned-original-source.tex, source lines 1896--1896. Navigation pointer invoked at proof line 1896; the counted span contains the pointer only and the cited result is not reproduced inside it, so the pointer is registered as a non-load-bearing navigation stub in the established form for Stacks material. No bibliography citation is involved.\par\medskip
\ifdefined\hypertarget\hypertarget{sourcepointer-3}{}\fi
\noindent\textit{Source reference.} Original lemma 010T: sources/stacks-sites-cohomology/pinned-original-source.tex, source lines 1900--1900. Navigation pointer invoked at proof line 1900; the counted span contains the pointer only and the cited result is not reproduced inside it, so the pointer is registered as a non-load-bearing navigation stub in the established form for Stacks material. No bibliography citation is involved.\par\medskip
\ifdefined\hypertarget\hypertarget{sourcepointer-4}{}\fi
\noindent\textit{Source reference.} Original lemma 03F9: sources/stacks-sites-cohomology/pinned-original-source.tex, source lines 1565--1565. Navigation pointer invoked at proof line 1904; the counted span contains the pointer only and the cited result is not reproduced inside it, so the pointer is registered as a non-load-bearing navigation stub in the established form for Stacks material. No bibliography citation is involved.\par\medskip
\ifdefined\hypertarget\hypertarget{sourcepointer-5}{}\fi
\noindent\textit{Source reference.} Original lemma 03FC: sources/stacks-sites-cohomology/pinned-original-source.tex, source lines 1816--1816. Navigation pointer invoked at proof line 1908; the counted span contains the pointer only and the cited result is not reproduced inside it, so the pointer is registered as a non-load-bearing navigation stub in the established form for Stacks material. No bibliography citation is involved.\par\medskip
\ifdefined\hypertarget\hypertarget{sourcepointer-6}{}\fi
\noindent\textit{Source reference.} Original lemma 03CI: sources/stacks-sites-cohomology/pinned-original-source.tex, source lines 1915--1915. Navigation pointer invoked at proof line 1915; the counted span contains the pointer only and the cited result is not reproduced inside it, so the pointer is registered as a non-load-bearing navigation stub in the established form for Stacks material. No bibliography citation is involved.\par\medskip
\begin{lemma}
\label{lemma-cohomology-modules-abelian-agree}
Let $\mathcal{C}$ be a site.
Let $\mathcal{O}$ be a sheaf of rings on $\mathcal{C}$.
Let $\mathcal{F}$ be an $\mathcal{O}$-module, and denote
$\mathcal{F}_{ab}$ the underlying sheaf of abelian groups.
Then we have
$$
H^i(\mathcal{C}, \mathcal{F}_{ab})
=
H^i(\mathcal{C}, \mathcal{F})
$$
and for any object $U$ of $\mathcal{C}$ we also have
$$
H^i(U, \mathcal{F}_{ab})
=
H^i(U, \mathcal{F}).
$$
Here the left hand side is cohomology computed in
$\textit{Ab}(\mathcal{C})$ and the right hand side
is cohomology computed in $\textit{Mod}(\mathcal{O})$.
\end{lemma}

\begin{proof}
By
Derived Categories, Lemma \ref{derived-lemma-higher-derived-functors}
the $\delta$-functor $(\mathcal{F} \mapsto H^p(U, \mathcal{F}))_{p \geq 0}$
is universal. The functor
$\textit{Mod}(\mathcal{O}) \to \textit{Ab}(\mathcal{C})$,
$\mathcal{F} \mapsto \mathcal{F}_{ab}$ is exact. Hence
$(\mathcal{F} \mapsto H^p(U, \mathcal{F}_{ab}))_{p \geq 0}$
is a $\delta$-functor also. Suppose we show that
$(\mathcal{F} \mapsto H^p(U, \mathcal{F}_{ab}))_{p \geq 0}$
is also universal. This will imply the second statement of the lemma
by uniqueness of universal $\delta$-functors, see
Homology, Lemma \ref{homology-lemma-uniqueness-universal-delta-functor}.
Since $\textit{Mod}(\mathcal{O})$ has enough injectives,
it suffices to show that $H^i(U, \mathcal{I}_{ab}) = 0$
for any injective object $\mathcal{I}$ in $\textit{Mod}(\mathcal{O})$, see
Homology, Lemma \ref{homology-lemma-efface-implies-universal}.

\medskip\noindent
Let $\mathcal{I}$ be an injective object of $\textit{Mod}(\mathcal{O})$.
Apply Lemma \ref{lemma-cech-vanish-collection}
with $\mathcal{F} = \mathcal{I}$, $\mathcal{B} = \mathcal{C}$
and $\text{Cov} = \text{Cov}_\mathcal{C}$.
Assumption (3) of that lemma holds by
Lemma \ref{lemma-injective-module-trivial-cech}.
Hence we see that $H^i(U, \mathcal{I}_{ab}) = 0$
for every object $U$ of $\mathcal{C}$.

\medskip\noindent
If $\mathcal{C}$ has a final
object then this also implies the first equality. If not, then
according to Sites, Lemma \ref{sites-lemma-topos-good-site} we see that
the ringed topos $(\Sh(\mathcal{C}), \mathcal{O})$ is equivalent to a
ringed topos where the underlying site does have a final object.
Hence the lemma follows.
\end{proof}

\end{document}
